\documentclass[a4paper,14pt]{extarticle}
\usepackage[utf8]{inputenc}
\usepackage[T2A]{fontenc}
\usepackage[russian]{babel}
\usepackage{geometry}
\geometry{left=3cm}
\geometry{right=1cm}
\geometry{top=2cm}
\geometry{bottom=2cm}
\usepackage{setspace}
\onehalfspacing
\usepackage{titlesec}
\usepackage{indentfirst}
\usepackage{enumitem}
\usepackage{tocloft}

% Настройка оформления разделов
\titleformat{\section}[block]
{\normalfont\bfseries\centering}
{\thesection}{1em}{}

\titleformat{\subsection}[block]
{\normalfont\bfseries\raggedright}
{\thesubsection}{1em}{}

% Настройка оглавления
\renewcommand{\cftsecleader}{\cftdotfill{\cftdotsep}}

\begin{document}


\tableofcontents
\newpage

\section*{Введение}
\addcontentsline{toc}{section}{Введение}

\textbf{Что здесь должно быть:}
\begin{itemize}
\item Объяснение проблемы: обычные ИИ-боты в играх ведут себя непредсказуемо
\item Зачем это нужно: чтобы создать более реалистичных и последовательных персонажей
\item Что сделали: разработали систему, которая добавляет эмоции и память ИИ-персонажу
\item Как проверяли: протестировали на диалогах и сравнили с обычным ботом
\end{itemize}

\textbf{На какую литературу ссылаться:}
\begin{itemize}
\item \cite{1} - примеры проблем в AI Dungeon
\item \cite{6} - современные подходы к ИИ-ведущим
\item \cite{9} - вызовы в диалоговых системах
\end{itemize}

\section{Анализ проблем существующих диалоговых систем}

\textbf{Что здесь должно быть:}
\begin{itemize}
\item Описание как работают текущие системы (типа AI Dungeon)
\item Главные проблемы:
  \begin{itemize}
  \item Нет памяти - забывают что было раньше
  \item Нет характера - сегодня добрый, завтра злой
  \item Слишком всем согласны - разрешают любые безумные действия
  \end{itemize}
\item Примеры из литературы где эти проблемы описываются
\end{itemize}

\textbf{На какую литературу ссылаться:}
\begin{itemize}
\item \cite{1}, \cite{2}, \cite{3} - примеры систем с проблемами
\item \cite{8} - анализ проблем в сторителлинге
\item \cite{9} - систематизация вызовов в диалоговых системах
\end{itemize}

\section{Архитектура и алгоритмы эмоционально-когнитивного агента}

\subsection{Модель личности агента: система ролей и принципов}

\textbf{Что здесь должно быть:}
\begin{itemize}
\item Как устроен класс Role - хранит описание персонажа
\item Что такое core\_principles и guidelines - правила поведения
\item Как это превращается в инструкцию для ИИ (prompt)
\item Пример кода с датаклассами
\end{itemize}

\textbf{На какую литературу ссылаться:}
\begin{itemize}
\item \cite{6} - использование интенций и теории сознания
\item \cite{12} - методы промпт-инжиниринга
\end{itemize}

\subsection{Когнитивная архитектура: обработка состояний и принятие решений}

\textbf{Что здесь должно быть:}
\begin{itemize}
\item Роль класса Brain - главный "мозг" системы
\item Что такое MoralScheme - схемы поведения
\item Как работает переключение между схемами
\item Пример с TransitionRule
\end{itemize}

\textbf{На какую литературу ссылаться:}
\begin{itemize}
\item \cite{10} - теория моральных основ
\item \cite{7} - сравнение статических и агентных подходов
\end{itemize}

\subsection{Динамическая модель эмоциональных состояний}

\textbf{Что здесь должно быть:}
\begin{itemize}
\item Математика эмоций: appraisals и feelings
\item Формулы обновления состояний
\item Что такое r\_const и p\_const (константы скорости изменений)
\item Как измеряется "расстояние" между эмоциями (euc\_dist)
\end{itemize}

\textbf{На какую литературу ссылаться:}
\begin{itemize}
\item \cite{11} - когнитивная структура эмоций
\item Marsella et al. - вычислительные модели эмоций
\end{itemize}

\subsection{Механизм взаимодействия с языковыми моделями}

\textbf{Что здесь должно быть:}
\begin{itemize}
\item Как класс Agent работает с ИИ
\item Метод generate\_changed\_message - формирование запроса
\item Как эмоции встраиваются в диалог
\item Пример кода с формированием промпта
\end{itemize}

\textbf{На какую литературу ссылаться:}
\begin{itemize}
\item \cite{5} - LLM как ассистенты
\item \cite{12} - методы промпт-инжиниринга
\end{itemize}

\section{Реализация и тестирование системы}

\subsection{Технологический стек и ключевые компоненты}

\textbf{Что здесь должно быть:}
\begin{itemize}
\item Что использовали: Python, aiohttp, numpy
\item Основные классы: ExamleAgent, Brain, MoralScheme
\item Ключевые методы и их назначение
\item Примеры кода из вашей реализации
\end{itemize}

\textbf{На какую литературу ссылаться:}
\begin{itemize}
\item Официальная документация Python, NumPy
\item \cite{6}, \cite{7} - аналогичные реализации
\end{itemize}

\subsection{Методика тестирования и экспериментальные сценарии}

\textbf{Что здесь должно быть:}
\begin{itemize}
\item Как тестировали: какие диалоги использовали
\item Что проверяли: стабильность эмоций, память, характер
\item С чем сравнивали: работа вашего агента vs обычная LLM
\item Метрики оценки: какие числа измеряли
\end{itemize}

\textbf{На какую литературу ссылаться:}
\begin{itemize}
\item \cite{6}, \cite{7}, \cite{8} - методики тестирования из схожих работ
\end{itemize}

\subsection{Анализ результатов и эффективности подхода}

\textbf{Что здесь должно быть:}
\begin{itemize}
\item Насколько лучше ваш агент по сравнению с базовой версией
\item Примеры удачных и неудачных диалогов
\item Выводы по эффективности подхода
\end{itemize}

\textbf{На какую литературу ссылаться:}
\begin{itemize}
\item Результаты собственного эксперимента
\item Сравнение с теоретическими положениями из Главы 1
\end{itemize}

\section*{Заключение}
\addcontentsline{toc}{section}{Заключение}

\textbf{Что здесь должно быть:}
\begin{itemize}
\item Кратко что сделали и что получилось
\item Главные достижения: что работает лучше
\item Где можно использовать вашу систему
\item Что можно улучшить в будущем
\item Итоговый вывод: стоит ли использовать такой подход
\end{itemize}

\textbf{На какую литературу ссылаться:}
\begin{itemize}
\item Обобщение результатов работы
\item Перспективы развития на основе \cite{5}, \cite{6}, \cite{7}
\end{itemize}

\newpage

\begin{thebibliography}{9}
\addcontentsline{toc}{section}{Список литературы}

\bibitem{1}
AI Dungeon. URL: \texttt{https://play.aidungeon.com/}

\bibitem{2}
AI Dungeon-Chi (ранний доступ). URL: \texttt{https://aidungeon-chi.vercel.app/}

\bibitem{3}
Dungeon Master AI. URL: \texttt{https://miniapps.ai/dungeon-master-ai/}

\bibitem{4}
HyperWrite AI – Interactive Story / Assistant Tool. URL: \texttt{https://app.hyperwriteai.com/}

\bibitem{5}
Zhu A., Martin L. J., Head A., Callison-Burch C. CALYPSO: LLMs as Dungeon Masters' Assistants. arXiv preprint, 2023. DOI: 10.48550/arXiv.2308.07540.

\bibitem{6}
Zhou P., Zhu A., Hu J., Pujara J., Ren X., Callison-Burch C., Choi Y., Ammanabrolu P. I Cast Detect Thoughts: Learning to Converse and Guide with Intents and Theory-of-Mind in Dungeons and Dragons. arXiv preprint, 2022. DOI: 10.48550/arXiv.2212.10060.

\bibitem{7}
Jørgensen N. H., Tharmabalan S., Aslan I., Hansen N. B., Merritt T. Static Vs. Agentic Game Master AI for Facilitating Solo Role-Playing Experiences. arXiv preprint, 2025. DOI: 10.48550/arXiv.2502.19519.

\bibitem{8}
Santiago J. M., Parayno R. L., Deja J. A., Samson B. P. V. Rolling the Dice: Imagining Generative AI as a Dungeons \& Dragons Storytelling Companion. arXiv preprint, 2023. DOI: 10.48550/arXiv.2304.01860.

\bibitem{9}
Callison-Burch C., Tomar G. S., Martin L. J., Ippolito D., Bailis S., Reitter D. Dungeons and Dragons as a Dialog Challenge for Artificial Intelligence. In: Proceedings of the 2022 Conference on Empirical Methods in Natural Language Processing, 2022.

\bibitem{10}
Haidt J. The Righteous Mind: Why Good People Are Divided by Politics and Religion. Pantheon Books, 2012.

\bibitem{11}
Ortony A., Clore G. L., Collins A. The Cognitive Structure of Emotions. Cambridge University Press, 1988.

\bibitem{12}
Liu P., et al. Pre-train, Prompt, and Predict: A Systematic Survey of Prompting Methods in Natural Language Processing. ACM Computing Surveys, 2023.

\end{thebibliography}

\newpage

\section*{Приложение А. Листинг ключевых компонентов системы}
\addcontentsline{toc}{section}{Приложение А}

\begin{verbatim}
# core/agent.py - базовый класс агента
class Agent(ABC):
    def __init__(self, agent_id: str, role: Role, 
                 brain: Brain, interface: Interface):
        self.id = agent_id
        self.role = role
        self.brain = brain
        self.interface = interface
        self.messages: List[Dict[str, str]] = [
            {"role": "system", "content": f"{self.role.to_string()}"}
        ]
    
    async def generate_reply(self, user_text: str) -> Dict[str, Any]:
        intents = await self.analyze_intentions(user_text)
        emotions = await self.analyze_emotions(user_text)
        self.brain.update_vectors(np.array(intents))
        
        context = {"last_message": user_text, 
                  "intents": intents, "emotions": emotions}
        
        if self.brain.check_transition(self.cur_scheme, context):
            self.brain.switch_to(self.cur_scheme + 1)
            
        changed_message = self.generate_changed_message(user_text, context)
        messages = list(self.messages)
        messages.append({"role": "user", "content": changed_message})
        reply_text = await self.interface.extract_text(messages)
        
        return reply_text
\end{verbatim}

\end{document}